\begin{table}[htbp]
\centering
\caption{Summary Statistics}
\label{tab:summary}
\begin{tabular}{lcc}
\hline\hline
Variable & Mean & SD \\
\hline
Math proficiency (\%) & 48.1 & 24.3 \\
Predicted exposure (std.) & -0.000 & 1.000 \\
Facilities within 50km & 73.7 & 82.2 \\
Mean distance to facility (km) & 29.9 & 10.0 \\
County SO\textsubscript{2} (ppb) & 2.83 & 3.33 \\
County PM\textsubscript{2.5} ($\mu$g/m$^3$) & 7.27 & 2.08 \\
\hline
Observations & \multicolumn{2}{c}{17,074} \\
Schools & \multicolumn{2}{c}{6,061} \\
Years & \multicolumn{2}{c}{2013, 2015, 2017} \\
\hline\hline
\end{tabular}
\begin{tablenotes}
\small
\item \textit{Notes:} Sample restricted to U.S. public schools (grades 3--8) within 50km of at least one EPA-monitored major air-emitting facility. Math proficiency is the school-level percentage of students scoring at or above proficiency on state assessments (EdFacts). Predicted exposure is the Gaussian plume dispersion-model predicted ground-level concentration (arbitrary units) from all nearby facilities, based on plant-school geometry and ASOS wind patterns.
\end{tablenotes}
\end{table}
